\clearpage
\SourceLeaf{148}
\begin{center}
\textsc{Gauss} an \textsc{Schumacher}. Göttingen, 6. Juli 1840.
\end{center}

\noindent . . . . . . . . . . .

Sie wissen, dass ich selbst auf das von mir seit 1794 gebrauchte Verfahren, dem später der Name Méthode des moindres quarrés gegeben wurde, niemals großen Wert gelegt habe. Verstehen Sie mich recht. Ich meine damit nicht den großen Nutzen, den das Verfahren bringt. Der ist klar genug, aber danach beurteile ich die Dinge nicht. Ich legte vielmehr deshalb oder insofern wenig Wert darauf, als mir der Gedanke von Anfang an so natürlich erschien und so außerordentlich nahe lag, dass ich nicht im Geringsten daran zweifelte, viele Menschen, die mit Zahlenrechnungen zu tun hatten, müssten von selbst auf einen solchen Kunstgriff gekommen sein und ihn gebraucht haben, ohne es deswegen der Mühe wert zu halten, um eine so natürliche Sache viel Aufhebens zu machen. Vor allem fiel mir \textsc{Tobias Mayer} ein, und ich erinnere mich sehr bestimmt, dass ich \textls{oft}, wenn ich mit anderen über meine Methode sprach (wie es z.\,B.\ während meiner Studienzeit 1795--1798 tatsächlich vielfach geschah), geäußert habe, ich wolle die allergrößte Wette eingehen, dass \textsc{Tobias Mayer} bei seinen Rechnungen dieselbe Methode schon gebraucht habe. Jetzt weiß ich aus jenen Papieren, dass ich jene Wette \textls{verloren} hätte. In der Tat enthalten sie Eliminationen, z.\,B.\ von 3 unbekannten Größen aus 4 oder 5 Gleichungen, doch so, wie der gewöhnlichste Rechner sie ausführen würde, ohne jede Spur irgendeiner feineren Kunst.

\noindent . . . . . . . . . . . .

